\section{Related Work}
\label{app:related_work}

\subsection{Extended Connections to Existing Work} \label{app:related_work_ratio_clipping}

In this work, we identify a structural flaw in PPO's ratio-clipping mechanism within the LLM regime: it over-penalizes low-probability tokens and under-penalizes high-probability ones, thereby impairing training efficiency and stability. Our proposed DPPO addresses this issue by directly constraining the policy divergence. 
This methodology aligns with the insights of \citet{wang2019trust, wang2020truly}, who observed similar exploration issues and proposed adaptive clipping based on KL divergence in traditional RL settings. However, in the context of LLMs, computing the exact divergence is prohibitive due to the huge memory footprint. To overcome this, we propose a binary divergence approximation, which empirically captures most of the benefits (see \Cref{fig:main_topk}). Furthermore, as demonstrated in \Cref{sec:stability} and \Cref{sec:training_efficiency}, the challenges of training stability and efficiency are exacerbated in LLMs by their expansive vocabularies, because low-probability tokens form a non-trivial portion of the entire distribution due to the long-tailed nature (see \Cref{fig:moe_prob_ratio_tv}). Finally, the training-inference mismatch inherent to the LLM era introduces additional algorithmic complexities, as further detailed in \Cref{sec:stability}.

\subsection{Training-inference Mismatch} \label{app:related_work_mismatch}
Recent work has identified a key culprit for training instability: the \textit{training-inference mismatch} (${\trainerpi}_\theta \neq {\rolloutpi}_\theta$), where the policy distribution used for gradient computation ($\trainerpi_\theta$) diverges from the one used for data generation ($\rolloutpi_\theta$), even when using identical model parameters $\theta$ \citep{yao2025offpolicy, qi2025defeating, liu-li-2025, zheng2025stabilizing}. This discrepancy arises from numerical precision errors \citep{qi2025defeating} and subtle differences in implementation \citep{Team2025EveryAM, he2025nondeterminism}. As training progresses, this mismatch can be amplified if the RL algorithm cannot manage it appropriately, leading to catastrophic performance degradation~\citep{qi2025defeating, liu-li-2025}.

Existing efforts to mitigate this issue primarily focus on correcting biased gradients through importance sampling. Building on this principle, techniques such as Truncated Importance Sampling (TIS)~\citep{yao2025offpolicy, zheng2025stabilizing} and Masked Importance Sampling \citep{liu-li-2025, team2025every} have been introduced at both the token and sequence levels. However, as suggested by \citet{qi2025defeating}, these methods often fail to achieve a satisfactory balance between training efficiency and stability. In contrast, our DPPO algorithm significantly enhances both aspects compared to these existing approaches.

Another line of research attempts to resolve the mismatch issue through higher precision \citep{qi2025defeating} or rigorous engineering alignment \citep{Team2025EveryAM, he2025nondeterminism, zhang2025deterministic}. While promising, these methods face limited applicability. For instance, aligning implementation details often requires specific training engines or model architectures, hindering broad adoption. Furthermore, in low-precision settings optimized for high-speed training, we must tolerate a significant training-inference mismatch. In such scenarios, a robust and fast algorithm like DPPO remains essential. Finally, our algorithmic design is orthogonal to these engineering-level optimizations and can be combined with them to achieve even greater performance gains.

